\scanpage{27}
Le fait qu'il soit gradué est immédiat. Il contient le groupe des cycles lin. équiv. à $0$, comme il résulte aussitôt \unsure{des définitions|de définition}, et du fait que sur $T=\mathbb P^1$, tout $0$-cycle $\unsure{\mathfrak a|\alpha|\mathfrak{t}}$ de degré $0$ satisfait $S_T(\mathfrak a)=0$. \deleted{Ainsi $N(X)$ est bien un \unsure{ss-|s-}groupe de $X$, contenant}

Pour prouver (iii), notons $z$ un cycle de la classe envisagée dans $A^*(\deleted{\unsure{T}}Y)$, et considérons un $Z$ sur $Y\times T$ et un $\mathfrak a$ comme dans 1.1., donc $z=Z(\mathfrak a)$. Si nous changeons $Z,\mathfrak a$ dans leurs classes d'équivalence linéaire \inserted{en $Z',\mathfrak a'$}, $Z'(\mathfrak a)$ sera \unsure{linéair\textsuperscript{mt}|linéairement} équivalent à $z$, et le Pb est le $\overline{\mathrm m}$ pour $z$ et $z'$.

Il suffit maintenant de prouver \deleted{$\unsure{Z'}$} qu'il existe $Z',\mathfrak a'$ tels que $(f\times\mathrm{id}_T)^*(Z')=Z'_1$ et $Z'(\mathfrak a'),\deleted{\unsure{Z_1(\mathfrak a')|Z'_1(\mathfrak{t}')}}$ soient définis ;

\diagram{\begin{tikzcd}[ampersand replacement=\&,column sep=small,row sep=small]
X\arrow[rr] \& \& Y\\
X\times T\arrow[rr,leftrightarrow]\arrow[dr] \& \&Y\times T\arrow[dl]\\
\deleted{\unsure{T}} \& T \&
\end{tikzcd}}

en effet, en vertu de la théorie des intersections, on aura alors $f^*(Z'(\mathfrak a'))=\deleted{\unsure{f^*}}(Z_1(\mathfrak a')$, or le premier mb \deleted{\unsure{représente est dans la|représentant et dans la}} représente $f^*(z)$ dans $A^*(X)$, \deleted{le deuxième est}

% archive sheet number: 26
